\section*{Chapitre \ref{ch:g11:polarity-and-cohesion} --- Électronégativité, polarité et interactions entre molécules}

\begin{solution}{exo:g11:polarity-and-cohesion:1}
\ce{H-Cl}~: le chlore (3.16 contre 2.20). \ce{O-H}~: l'oxygène (3.44).
\ce{C-O}~: l'oxygène (3.44 contre 2.55).
\end{solution}

\begin{solution}{exo:g11:polarity-and-cohesion:2}
\ce{H-H}~: 0, non polarisée. \ce{C-H}~: 0.35, considérée comme non
polarisée. \ce{O-H}~: 1.24, polarisée. \ce{N-H}~: 0.84, polarisée.
\ce{Cl-Cl}~: 0, non polarisée. \ce{C-Cl}~: 0.61, polarisée.
\end{solution}

\begin{solution}{exo:g11:polarity-and-cohesion:3}
L'azote (3.04 > 2.19) et l'oxygène (3.44 > 2.58)~: l'\omterm{def:g11:polarity-and-cohesion:electronegativity}{électronégativité}
diminue en descendant dans un groupe. L'azote (3.04 > 2.55) et le
magnésium (1.31 > 0.93)~: elle augmente de gauche à droite le long
d'une \omterm{def:g9:periodic-table-first-look:period}{période}.
\end{solution}

\begin{solution}{exo:g11:polarity-and-cohesion:4}
L'ammoniac et l'eau~: chacun possède des \omterm{def:g7:atoms-and-molecules:atom}{atomes} d'hydrogène liés à
l'azote ou à l'oxygène, et des \omterm{def:g10:lewis-and-shape:lone-pair}{doublets non liants} sur cet \omterm{def:g7:atoms-and-molecules:atom}{atome}. Le
méthane n'a pas d'hydrogène lié à N, O ou F~; le chlorure d'hydrogène a
un hydrogène lié au chlore, qui ne fait pas partie des trois \omterm{def:g7:atoms-and-molecules:atom}{atomes} de
la définition.
\end{solution}

\begin{solution}{exo:g11:polarity-and-cohesion:5}
Oui~: des \omterm{def:g11:polarity-and-cohesion:van-der-waals}{interactions de van der Waals} existent entre toutes les
\omterm{def:g7:atoms-and-molecules:molecule}{molécules}. Mais les \omterm{def:g7:atoms-and-molecules:molecule}{molécules} de méthane sont petites (10 \omterm{def:g9:inside-the-atom:nucleus}{électrons})~:
ces interactions sont très faibles et se rompent à très basse
température~; le méthane bout vers \qty{-162}{\celsius}.
\end{solution}

\begin{solution}{exo:g11:polarity-and-cohesion:6}
\ce{CH2Cl2}~: polaire (deux liaisons \ce{C-Cl} polarisées et deux
liaisons \ce{C-H} presque non polarisées~: aucune symétrie ne les
compense). \ce{CCl4}~: apolaire (quatre liaisons identiques aux sommets
d'un tétraèdre). \ce{NH3}~: polaire (pyramide). \ce{CO2}~: apolaire
(linéaire, liaisons opposées). \ce{BF3}~: apolaire (trois \omterm{def:g11:polarity-and-cohesion:polar-bond}{liaisons
polarisées} identiques à \ang{120} dans un plan se compensent).
\end{solution}

\begin{solution}{exo:g11:polarity-and-cohesion:7}
L'hydrogène du groupe \ce{O-H} peut être donné (pas ceux de \ce{CH3},
liés au carbone)~; l'\omterm{def:g7:atoms-and-molecules:atom}{atome} d'oxygène, avec ses deux \omterm{def:g10:lewis-and-shape:lone-pair}{doublets non liants},
peut en recevoir. Oui~: le \ce{O-H} du méthanol peut se lier à
l'oxygène d'une \omterm{def:g7:atoms-and-molecules:molecule}{molécule} d'eau, et un \ce{O-H} de l'eau à l'oxygène du
méthanol, par exemple \ce{CH3-O-H}\,$\cdots$\,\ce{OH2}.
\end{solution}

\begin{solution}{exo:g11:polarity-and-cohesion:8}
Les trois \omterm{def:g7:atoms-and-molecules:molecule}{molécules} sont apolaires et de même géométrie, mais de plus
en plus grosses (10, 18 et 36 \omterm{def:g9:inside-the-atom:nucleus}{électrons})~: leurs \omterm{def:g11:polarity-and-cohesion:van-der-waals}{interactions de van der
Waals} augmentent, et leurs températures d'ébullition aussi.
\end{solution}

\begin{solution}{exo:g11:polarity-and-cohesion:9}
Le \omterm{def:g9:periodic-table-first-look:period}{groupe} 14~: le méthane n'a pas d'hydrogène lié à N, O ou F et ne
forme pas de \omterm{def:g11:polarity-and-cohesion:hydrogen-bond}{liaisons hydrogène}~; il suit donc la tendance de son
groupe.
\end{solution}

\begin{solution}{exo:g11:polarity-and-cohesion:10}
Les \omterm{def:g11:polarity-and-cohesion:van-der-waals}{interactions de van der Waals}~: les \omterm{def:g7:atoms-and-molecules:molecule}{molécules} grossissent de
\ce{HCl} à \ce{HI}, et l'augmentation de ces interactions avec la taille
l'emporte sur la diminution de l'attraction entre les \omterm{def:g11:polarity-and-cohesion:polar-bond}{charges
partielles}.
\end{solution}

\begin{solution}{exo:g11:polarity-and-cohesion:11}
Une \omterm{def:g7:atoms-and-molecules:molecule}{molécule} de diiode est bien plus grosse qu'une \omterm{def:g7:atoms-and-molecules:molecule}{molécule} de dichlore
(106 \omterm{def:g9:inside-the-atom:nucleus}{électrons} contre 34)~: ses \omterm{def:g11:polarity-and-cohesion:van-der-waals}{interactions de van der Waals} sont bien
plus fortes, assez pour maintenir les \omterm{def:g7:atoms-and-molecules:molecule}{molécules} dans un solide à
température ambiante.
\end{solution}

\begin{solution}{exo:g11:polarity-and-cohesion:12}
Le propane est apolaire~: seulement des \omterm{def:g11:polarity-and-cohesion:van-der-waals}{interactions de van der Waals}.
Le méthoxyméthane est polaire (deux liaisons \ce{C-O} polarisées dans un
enchaînement coudé \ce{C-O-C}) mais n'a pas d'hydrogène sur son
oxygène~: pas de \omterm{def:g11:polarity-and-cohesion:hydrogen-bond}{liaisons hydrogène} entre ses \omterm{def:g7:atoms-and-molecules:molecule}{molécules}. L'éthanol est
polaire et ses groupes \ce{O-H} forment des \omterm{def:g11:polarity-and-cohesion:hydrogen-bond}{liaisons hydrogène}. Plus les
attractions sont fortes, plus la température d'ébullition est élevée~:
propane < méthoxyméthane < éthanol.
\end{solution}

\begin{solution}{exo:g11:polarity-and-cohesion:13}
Le réseau ouvert de la glace occupe plus de place que les mêmes
\omterm{def:g7:atoms-and-molecules:molecule}{molécules} dans le liquide~: la glace est moins dense que l'eau liquide
et flotte. Un lac gèle par le haut~: la \omterm{def:g10:electron-shells:configuration}{couche} de glace isole l'eau qui
se trouve dessous, qui reste liquide, et les poissons survivent.
\end{solution}

\begin{solution}{exo:g11:polarity-and-cohesion:14}
Le méthane solide (\omterm{def:g11:polarity-and-cohesion:van-der-waals}{interactions de van der Waals} entre petites
\omterm{def:g7:atoms-and-molecules:molecule}{molécules}) < la glace (\omterm{def:g11:polarity-and-cohesion:hydrogen-bond}{liaisons hydrogène}) < le chlorure de potassium
(attraction entre \omterm{def:g9:ions:ion}{ions}).
\end{solution}

\begin{solution}{exo:g11:polarity-and-cohesion:15}
Une \omterm{def:g7:atoms-and-molecules:molecule}{molécule} d'eau a deux \omterm{def:g7:atoms-and-molecules:atom}{atomes} d'hydrogène à donner et deux \omterm{def:g10:lewis-and-shape:lone-pair}{doublets
non liants} pour en recevoir~: chaque \omterm{def:g7:atoms-and-molecules:molecule}{molécule} peut participer à quatre
\omterm{def:g11:polarity-and-cohesion:hydrogen-bond}{liaisons hydrogène}, deux données et deux reçues, ce qui construit un
réseau complet. Une \omterm{def:g7:atoms-and-molecules:molecule}{molécule} de fluorure d'hydrogène a trois \omterm{def:g10:lewis-and-shape:lone-pair}{doublets
non liants} mais un seul \omterm{def:g7:atoms-and-molecules:atom}{atome} d'hydrogène~: elle ne donne qu'une \omterm{def:g11:polarity-and-cohesion:hydrogen-bond}{liaison
hydrogène}, si bien qu'il ne se forme en moyenne qu'une liaison par
\omterm{def:g7:atoms-and-molecules:molecule}{molécule} (chaînes en zigzag). L'eau a environ deux fois plus de \omterm{def:g11:polarity-and-cohesion:hydrogen-bond}{liaisons
hydrogène} par \omterm{def:g7:atoms-and-molecules:molecule}{molécule}, et bout plus haut.
\end{solution}

\begin{solution}{pb:g11:polarity-and-cohesion:1}
\textbf{1.} \ce{H-O-H}, \ce{H-S-H}, \ce{H-Se-H}, chaque \omterm{def:g7:atoms-and-molecules:atom}{atome} central
portant deux \omterm{def:g10:lewis-and-shape:lone-pair}{doublets non liants}~: O, S et Se ont six \omterm{def:g10:electron-shells:valence}{électrons de
valence}, car ils sont dans le même groupe.

\textbf{2.} $M(\ce{H2O}) = \qty{18.0}{g/mol}$, $M(\ce{H2S}) =
\qty{34.1}{g/mol}$, $M(\ce{H2Se}) = \qty{81.0}{g/mol}$.

\textbf{3.} $3.44 - 2.20 = 1.24$~; $2.58 - 2.20 = 0.38$~;
$2.55 - 2.20 = 0.35$. Seule la liaison \ce{O-H} est nettement
polarisée.

\textbf{4.} Coudée (quatre groupes, dont deux \omterm{def:g10:lewis-and-shape:lone-pair}{doublets non liants}). L'eau
est polaire~; le sulfure et le séléniure d'hydrogène ne le sont que très
peu.

\textbf{5.} 10, 18 et 36 \omterm{def:g9:inside-the-atom:nucleus}{électrons}~: c'est le séléniure d'hydrogène qui
a les \omterm{def:g11:polarity-and-cohesion:van-der-waals}{interactions de van der Waals} les plus fortes.

\textbf{6.} L'eau seulement~: ses \omterm{def:g7:atoms-and-molecules:atom}{atomes} d'hydrogène sont liés à
l'oxygène, l'un des trois \omterm{def:g7:atoms-and-molecules:atom}{atomes} très électronégatifs~; le soufre et le
sélénium ne sont pas assez électronégatifs pour donner à leurs \omterm{def:g7:atoms-and-molecules:atom}{atomes}
d'hydrogène un $\delta^+$ marqué.

\textbf{7.} Quatre~: deux par ses propres \omterm{def:g7:atoms-and-molecules:atom}{atomes} d'hydrogène, deux
reçues par ses deux \omterm{def:g10:lewis-and-shape:lone-pair}{doublets non liants}.

\textbf{8.} Des \omterm{def:g11:polarity-and-cohesion:van-der-waals}{interactions de van der Waals} (y compris une faible
attraction entre les petites \omterm{def:g11:polarity-and-cohesion:polar-bond}{charges partielles}).

\textbf{9.} $212.9 - 273.15 = \qty{-60.3}{\celsius}$ et
\qty{-41.3}{\celsius}.

\textbf{10.} $-41.3 - (-60.3) = \qty{19.0}{\celsius}$.

\textbf{11.} La \omterm{def:g7:atoms-and-molecules:molecule}{molécule} est plus grosse~: ses \omterm{def:g11:polarity-and-cohesion:van-der-waals}{interactions de van der
Waals} sont plus fortes.

\textbf{12.} $-60.3 - 19.0 = \qty{-79.3}{\celsius}$.

\textbf{13.} Augmentation~: $-88.1 - (-112) = \qty{23.9}{\celsius}$~;
prévision pour le méthane~: $-112 - 23.9 = \qty{-136}{\celsius}$, contre
\qty{-162}{\celsius} en réalité~: l'extrapolation est trop haute
d'environ \qty{26}{\celsius}. Elle n'est pas exacte~; la tendance réelle
n'est pas une droite.

\textbf{14.} Augmentation~: $\qty{18.7}{\celsius}$~; prévision pour le
fluorure d'hydrogène~: $-85.1 - 18.7 = \qty{-103.8}{\celsius}$, contre
\qty{19.6}{\celsius}~: un écart d'environ \qty{123}{\celsius}.

\textbf{15.} $100 - (-79.3) = \qty{179}{\celsius}$, environ
\qty{180}{\celsius}.

\textbf{16.} Augmentation~: $\qty{25.3}{\celsius}$~; prévision~:
$-87.8 - 25.3 =
\qty{-113}{\celsius}$~; écart~: $-33.4 - (-113) = \qty{80}{\celsius}$.

\textbf{17.} L'eau (\qty{180}{\celsius}) > le fluorure d'hydrogène
(\qty{123}{\celsius}) > l'ammoniac (\qty{80}{\celsius}). L'eau donne
deux \omterm{def:g7:atoms-and-molecules:atom}{atomes} d'hydrogène et en reçoit avec deux \omterm{def:g10:lewis-and-shape:lone-pair}{doublets non liants}~:
quatre \omterm{def:g11:polarity-and-cohesion:hydrogen-bond}{liaisons hydrogène} par \omterm{def:g7:atoms-and-molecules:molecule}{molécule}, toutes utilisées. Le fluorure
d'hydrogène n'a qu'un hydrogène à donner, et l'ammoniac qu'un \omterm{def:g10:lewis-and-shape:lone-pair}{doublet
non liant} pour recevoir~: en moyenne moins de liaisons par \omterm{def:g7:atoms-and-molecules:molecule}{molécule}, et
l'azote, moins électronégatif, en forme de plus faibles.

\textbf{18.} Gazeuse~: elle bouillirait vers \qty{-80}{\celsius}. La
Terre n'aurait pas d'eau liquide, pas d'océans, pas de pluie, et aucune
des formes de vie qui en dépendent.

\textbf{19.} On extrapole une droite tracée par deux points seulement,
et le test sur le \omterm{def:g9:periodic-table-first-look:period}{groupe} 14 montre qu'une telle extrapolation peut se
tromper de quelque \qty{25}{\celsius}.

\textbf{20.} Sans \omterm{def:g11:polarity-and-cohesion:hydrogen-bond}{liaisons hydrogène}, l'eau bouillirait vers
\qty{-80}{\celsius}, quelque \qty{180}{\celsius} au-dessous de sa
véritable température d'ébullition.
\end{solution}
